\documentclass[11pt,a4paper]{article}
\usepackage[T1]{fontenc}
\usepackage{amsmath}
\usepackage{amssymb}
\usepackage[table]{xcolor}
\usepackage{array}
\usepackage{tabularx}
\usepackage{multirow}
\usepackage{geometry}
\usepackage{enumitem}
\usepackage{parskip}
\usepackage{titlesec}
\usepackage{tocloft}
\usepackage{needspace}
\usepackage[most]{tcolorbox}

\geometry{margin=2.2cm}

\titleformat{\section}{\large\bfseries}{}{0em}{}
\setlength{\cftbeforesecskip}{6pt}

% Question heading that refuses to be orphaned at the foot of a page.
% Reserves 8 lines so the heading stays with the Mark Scheme table beneath it.
\newcommand{\qhead}[2][]{\needspace{8\baselineskip}\section[#1]{#2}}

% Operators and shorthands used across 4037
\DeclareMathOperator{\cosec}{cosec}
\newcommand{\dd}{\mathrm{d}}
\newcommand{\dydx}{\dfrac{\dd y}{\dd x}}
\newcommand{\ncr}[2]{{}^{#1}\mathrm{C}_{#2}}
\newcommand{\npr}[2]{{}^{#1}\mathrm{P}_{#2}}

% Pastel colours, cycled across the questions (q1..q9)
\definecolor{q1color}{HTML}{D6EAF8} % pastel blue
\definecolor{q2color}{HTML}{D5F5E3} % pastel green
\definecolor{q3color}{HTML}{FDEBD0} % pastel orange
\definecolor{q4color}{HTML}{F9E79F} % pastel yellow
\definecolor{q5color}{HTML}{E8DAEF} % pastel purple
\definecolor{q6color}{HTML}{FADBD8} % pastel red/pink
\definecolor{q7color}{HTML}{D1F2EB} % pastel teal
\definecolor{q8color}{HTML}{FCF3CF} % pastel cream
\definecolor{q9color}{HTML}{EBDEF0} % pastel lavender

% Mark-scheme table (FOUR columns: Question | Answer | Marks | Partial Marks or Guidance)
% Optional argument sets the fourth heading; default is Partial Marks.
\newenvironment{mstab}[2][Partial Marks]{%
  {\color{#2!75!black}\nobreak Mark Scheme}\par\nopagebreak\smallskip
  \arrayrulecolor{#2!70!black}%
  \renewcommand{\arraystretch}{1.35}%
  \tabularx{\textwidth}{%
    |>{\columncolor{#2!12}\raggedright\arraybackslash}p{1.6cm}|%
     >{\columncolor{#2!12}\raggedright\arraybackslash}p{5.2cm}|%
     >{\columncolor{#2!12}\centering\arraybackslash}p{1.2cm}|%
     >{\columncolor{#2!12}\raggedright\arraybackslash}X|}%
  \hline
  \rowcolor{#2!45} Question & Answer & Marks & #1 \\ \hline
}{%
  \endtabularx
  \arrayrulecolor{black}%
  \par\medskip
}

% Full-width banner row for Alternative methods: \altrow{q3color}{Alternative 1}
\newcommand{\altrow}[2]{\multicolumn{4}{|l|}{\cellcolor{#1!30}\textbf{#2}} \\ \hline}

% Exemplar-answer box
\newtcolorbox{ansbox}[1]{
  colback=#1!8,
  colbacktitle=#1!30,
  colframe=#1!70!black,
  coltitle=black,
  fonttitle=\normalfont,
  title=Worked Solution,
  boxrule=0.6pt,
  arc=2mm,
  left=4pt, right=4pt, top=4pt, bottom=4pt,
  breakable
}

\newcommand{\topiclabel}[1]{\hfill\normalfont\small\itshape\textcolor{black!55}{#1}}

% Per-question head: \examq{paper-id}{number}{topics}{marks}
% Same command and rendered line as the question paper.
\newcommand{\examq}[4]{\section[#1 - Q#2]{#1 - Q#2 - #3 - #4}}

\begin{document}

\begin{center}
  {\LARGE \textbf{Cambridge O Level Additional Mathematics 4037/11}}\\[6pt]
  {\large May/June 2025 - Paper 1 (Non-calculator): Mark Scheme \& Worked Solutions}
\end{center}

\bigskip

%%%%%%%%%%%%%%%%%%%
\examq{4037/11/M/J/25}{1}{Vectors in Two Dimensions}{5}
%%%%%%%%%%%%%%%%%%%
\begin{mstab}{q1color}
1(a) & $\begin{pmatrix} -2.5 \\ 5 \end{pmatrix}$ oe & \textbf{B2} & \textbf{B1} for $-\frac{1}{4}\begin{pmatrix} -2 \\ 8 \end{pmatrix} \left[ + \begin{pmatrix} -3 \\ 7 \end{pmatrix} \right]$ oe, soi \\ \hline
\multirow{2}{*}{1(b)} & $3\alpha - 8 = 2\alpha + 5\beta$ \newline or $18 - 2\beta = 20$ & \textbf{M1} & Equates like vectors at least once \\ \cline{2-4}
 & $\alpha = 3$, $\beta = -1$ & \textbf{A2} & \textbf{A1} for either value correct \\ \hline
\end{mstab}

\begin{ansbox}{q1color}
(a) Finding vector $\overrightarrow{RQ}$\\[4pt]
We are given $\overrightarrow{PQ} = \begin{pmatrix} -3 \\ 7 \end{pmatrix}$ and $4\overrightarrow{PR} = \begin{pmatrix} -2 \\ 8 \end{pmatrix}$. First, find $\overrightarrow{PR}$:
\[
\overrightarrow{PR} = \frac{1}{4}\begin{pmatrix} -2 \\ 8 \end{pmatrix} = \begin{pmatrix} -0.5 \\ 2 \end{pmatrix}
\]
Using vector addition, $\overrightarrow{PR} + \overrightarrow{RQ} = \overrightarrow{PQ}$, so:
\[
\overrightarrow{RQ} = \overrightarrow{PQ} - \overrightarrow{PR} = \begin{pmatrix} -3 \\ 7 \end{pmatrix} - \begin{pmatrix} -0.5 \\ 2 \end{pmatrix} = \boxed{\begin{pmatrix} -2.5 \\ 5 \end{pmatrix}}
\]

(b) Finding the values of $\alpha$ and $\beta$\\[4pt]
Given vectors $\mathbf{a} = \alpha\mathbf{i} + 6\mathbf{j}$, $\mathbf{b} = 4\mathbf{i} + \beta\mathbf{j}$, and $\mathbf{c} = (2\alpha + 5\beta)\mathbf{i} + 20\mathbf{j}$. We are given that $\mathbf{c} = 3\mathbf{a} - 2\mathbf{b}$. Evaluating the right-hand side:
\[
3\mathbf{a} - 2\mathbf{b} = 3(\alpha\mathbf{i} + 6\mathbf{j}) - 2(4\mathbf{i} + \beta\mathbf{j}) = (3\alpha - 8)\mathbf{i} + (18 - 2\beta)\mathbf{j}
\]
Equating components with $\mathbf{c}$:
\[
\text{\textbf{i} components: } 2\alpha + 5\beta = 3\alpha - 8 \implies \alpha - 5\beta = 8
\]
\[
\text{\textbf{j} components: } 20 = 18 - 2\beta \implies 2\beta = -2 \implies \beta = -1
\]
Substitute $\beta = -1$ into the $\mathbf{i}$ component equation:
\[
\alpha - 5(-1) = 8 \implies \alpha + 5 = 8 \implies \alpha = 3
\]
Thus, $\boxed{\alpha = 3, \beta = -1}$.
\end{ansbox}

%%%%%%%%%%%%%%%%%%%
\examq{4037/11/M/J/25}{2}{Quadratic Functions}{4}
%%%%%%%%%%%%%%%%%%%
\begin{mstab}{q2color}
2 & $5x^2 - 10x - 15 \ge 0$ oe & \textbf{M1} & where $*$ is any inequality sign or = \\ \cline{2-4}
 & Factorises their 3-term quadratic expression or solves their 3-term quadratic equation & \textbf{M1} & \\ \cline{2-4}
 & Critical values $-1$ and $3$ & \textbf{A1} & \\ \cline{2-4}
 & $-1 \le x \le 3$ & \textbf{A1} & Do not accept separate inequalities unless connected with \textquotedbl{}and\textquotedbl{} \\ \hline
\end{mstab}

\begin{ansbox}{q2color}
Solving the inequality $(3-x)(5x+8) \ge 9-3x$\\[4pt]
Expand the left-hand side and simplify the inequality:
\[
15x + 24 - 5x^2 - 8x \ge 9 - 3x
\]
\[
-5x^2 + 7x + 24 \ge 9 - 3x
\]
Rearranging all terms to one side to form a quadratic expression $\ge 0$:
\[
-5x^2 + 10x + 15 \ge 0 \implies 5x^2 - 10x - 15 \le 0
\]
Dividing through by $5$:
\[
x^2 - 2x - 3 \le 0
\]
Factorising the quadratic:
\[
(x-3)(x+1) \le 0
\]
The critical values are $x = -1$ and $x = 3$. For the expression to be less than or equal to zero, $x$ must lie between the roots:
\[
\boxed{-1 \le x \le 3}
\]
\end{ansbox}

%%%%%%%%%%%%%%%%%%%
\examq{4037/11/M/J/25}{3}{Coordinate Geometry of the Circle}{6}
%%%%%%%%%%%%%%%%%%%
\begin{mstab}{q3color}
3(a) & $(3-4)^2 + (-1+3)^2 = 5$ \newline or showing that distance between point $A$ and centre of circle $=$ radius e.g. $\sqrt{(3-4)^2 + (-1+3)^2} = \sqrt{5}$ & \textbf{B1} & Accept if $x$-coordinate substituted to find $y$-coordinate \\ \hline
\multirow{2}{*}{3(b)} & $\begin{pmatrix} 4 \\ -3 \end{pmatrix} + \begin{pmatrix} 1 \\ -2 \end{pmatrix}$ oe, soi & \textbf{M1} & Award \textbf{M1} for the 2 equations or using centre and point $A$ e.g. $\frac{3+x}{2} = 4$ and $\frac{-1+y}{2} = -3$ \\ \cline{2-4}
 & $(5,\,-5)$ & \textbf{A1} & \\ \hline
3(c) & gradient of radius $=\frac{-3+1}{4-3} = -2$ soi by gradient of tangent & \textbf{M1} & \\ \cline{2-4}
 & gradient of tangent $=\frac{-1}{\text{their }(-2)}$ & \textbf{M1} & \textbf{FT} $\text{their } -2$ \\ \cline{2-4}
 & $y+1 = \frac{1}{2}(x-3)$ oe & \textbf{A1} & ISW from a correct unsimplified answer \\ \hline
\altrow{q3color}{Alternative}
 & use of differentiation e.g.: \newline $(y+3)^2 = 5 - (x-4)^2$ \newline $y = \sqrt{-x^2+8x-11} - 3$ \newline $\dfrac{\dd y}{\dd x} = \dfrac{-2x+8}{2\sqrt{-x^2+8x-11}}$ \newline Substitute $(3,-1)$ in $\text{their } \dfrac{\dd y}{\dd x}$ \newline $= \frac{1}{2}$ to get gradient & \textbf{(M1)} & allow one error \\ \hline
 & or use of implicit differentiation (not on syllabus) e.g.: \newline $2x + 2y\dfrac{\dd y}{\dd x} - 8 + 6\dfrac{\dd y}{\dd x} = 0$ \newline leading to $\dfrac{\dd y}{\dd x} = \dfrac{8-2x}{2y+6}$ & \textbf{(M1)} & Dep on first \textbf{M1} \\ \hline
 & $y+1 = \frac{1}{2}(x-3)$ oe & \textbf{(A1)} & ISW from a correct unsimplified answer \\ \hline
\end{mstab}

\begin{ansbox}{q3color}
(a) Showing $A$ lies on the circle\\[4pt]
The circle has equation $(x-4)^2 + (y+3)^2 = 5$ with centre $(4,-3)$ and radius $\sqrt{5}$. Substitute point $A(3,-1)$ into the left-hand side:
\[
(3-4)^2 + (-1+3)^2 = (-1)^2 + 2^2 = 1 + 4 = 5
\]
Since this matches the right-hand side, point $A$ lies on the circumference \quad \text{as required.}

(b) Coordinates of $B$\\[4pt]
Since $AB$ is a diameter, the centre $(4,-3)$ is the midpoint of $A(3,-1)$ and $B(x, y)$. Using midpoint formulas:
\[
\frac{3+x}{2} = 4 \implies 3+x = 8 \implies x = 5
\]
\[
\frac{-1+y}{2} = -3 \implies -1+y = -6 \implies y = -5
\]
Thus, the coordinates of $B$ are $\boxed{(5,\,-5)}$.

(c) Equation of the tangent at $A$\\[4pt]
The radius connects the centre $(4,-3)$ to $A(3,-1)$. Its gradient is:
\[
m_{\text{radius}} = \frac{-1 - (-3)}{3 - 4} = \frac{2}{-1} = -2
\]
The tangent is perpendicular to the radius, so its gradient is the negative reciprocal:
\[
m_{\text{tangent}} = -\frac{1}{-2} = \frac{1}{2}
\]
Using the point-slope form with point $A(3,-1)$:
\[
y - (-1) = \frac{1}{2}(x - 3) \implies \boxed{y + 1 = \frac{1}{2}(x - 3)}
\]
\end{ansbox}

%%%%%%%%%%%%%%%%%%%
\examq{4037/11/M/J/25}{4}{Quadratic Functions \textperiodcentered\ Simultaneous Equations \textperiodcentered\ Logarithmic \& Exponential Functions}{9}
%%%%%%%%%%%%%%%%%%%
\begin{mstab}{q4color}
4(a) & Solves or factorises $x^{\frac{1}{3}} - x^{\frac{1}{6}} - 2 = 0$: \newline $\left(x^{\frac{1}{6}} + 1\right)\left(x^{\frac{1}{6}} - 2\right)$ \newline or when substituting $y = x^{\frac{1}{6}}$ factorises to $\left(y+1\right)\left(y-2\right)$ & \textbf{M1} & A substitution may be used. Awarded for $(x^{\frac{1}{3}})^6 - (x^{\frac{1}{6}})^6 = 2^6$ scores 0 marks. $(x^{\frac{1}{3}}) - (x^{\frac{1}{6}}) = x^{\log_x 2}$ then $\frac{1}{3} - \frac{1}{6} = \log_x 2$ scores 0 marks. \\ \cline{2-4}
 & $\left[x^{\frac{1}{6}} = -1\right]$, $x^{\frac{1}{6}} = 2$ & \textbf{A1} & \\ \cline{2-4}
 & $x = 64$ \newline $\left[x = 1\right]$ & \textbf{A1} & \\ \cline{2-4}
 & Rejects $x^{\frac{1}{6}} = -1$ or $x=1$ ignored at some stage, soi & \textbf{B1} & \\ \hline
\multirow{2}{*}{4(b)} & $x + 2y = 1$ \newline Correctly eliminates one unknown e.g. $x = 1 - 2y$ or $y = \frac{1-x^2}{4x+1}$ & \textbf{B1} \newline \textbf{M1} & Dep on \textbf{B1}, allow unsimplified \\ \cline{2-4}
 & Correct quadratic in solvable form: $y - 4y^2 [= 0]$ oe or $2x^2 - 3x + 1 [= 0]$ oe & \textbf{A1} & \\ \cline{2-4}
 & Factorises or solves their quadratic to get two solutions & \textbf{M1} & Dep on previous \textbf{M1} \\ \cline{2-4}
 & $y = \frac{1}{4}, x = \frac{1}{2}$ and $y = 0, x = 1$ & \textbf{A1} & \\ \hline
\end{mstab}

\begin{ansbox}{q4color}
(a) Solving the equation $x^{\frac{1}{3}} - x^{\frac{1}{6}} = 2$\\[4pt]
Let $y = x^{\frac{1}{6}}$, so $y^2 = x^{\frac{1}{3}}$. Substituting into the equation:
\[
y^2 - y = 2 \implies y^2 - y - 2 = 0
\]
Factorising the quadratic:
\[
(y - 2)(y + 1) = 0 \implies y = 2 \text{ or } y = -1
\]
Substitute back $x^{\frac{1}{6}} = y$:
* Case 1: $x^{\frac{1}{6}} = 2 \implies x = 2^6 = 64$.
* Case 2: $x^{\frac{1}{6}} = -1$, which has no real solution since principal even roots cannot be negative.

Thus, the only valid solution is $\boxed{x = 64}$.

(b) Solving simultaneous equations\\[4pt]
Given equations:
1) $\lg(x + 2y) = 0$
2) $x^2 + 4xy + y = 1$

From equation (1):
\[
\lg(x + 2y) = 0 \implies x + 2y = 10^0 = 1 \implies x = 1 - 2y
\]
Substitute $x = 1 - 2y$ into equation (2):
\[
(1 - 2y)^2 + 4(1 - 2y)y + y = 1
\]
\[
1 - 4y + 4y^2 + 4y - 8y^2 + y = 1
\]
Simplify terms:
\[
-4y^2 + y = 0 \implies 4y^2 - y = 0 \implies y(4y - 1) = 0
\]
So $y = 0$ or $y = \frac{1}{4}$.
* If $y = 0$: $x = 1 - 2(0) = 1$.
* If $y = \frac{1}{4}$: $x = 1 - 2\left(\frac{1}{4}\right) = 1 - \frac{1}{2} = \frac{1}{2}$.

Thus, the solutions are $\boxed{x = 1, y = 0 \text{ and } x = \frac{1}{2}, y = \frac{1}{4}}$.
\end{ansbox}

%%%%%%%%%%%%%%%%%%%
\examq{4037/11/M/J/25}{5}{Equations, Inequalities \& Graphs \textperiodcentered\ Trigonometry}{5}
%%%%%%%%%%%%%%%%%%%
\begin{mstab}{q5color}
5(a) & $2x - 1 = -3$ and $2x - 1 = 3$ & \textbf{M1} & \\ \cline{2-4}
 & $x = 2$, $x = -1$ & \textbf{A2} & \textbf{A1} for either correct \\ \hline
\altrow{q5color}{Alternative}
 & $4x^2 - 4x - 8 = 0$ oe & \textbf{(B1)} & \\ \hline
 & Solves or factorises their 3-term quadratic & \textbf{(M1)} & \\ \hline
 & $x = -1, 2$ & \textbf{(A1)} & \\ \hline
\multirow{2}{*}{5(b)} & Correct graph (one full sine cycle, maximum $3$, minimum $-7$) & \textbf{B1} & \\ \cline{2-4}
 & Correct graph with midline $y = -2$ & \textbf{B1} & \\ \hline
\end{mstab}

\begin{ansbox}{q5color}
(a) Solving the equation $5\lvert 2x - 1 \rvert + 8 = 23$\\[4pt]
Isolate the modulus term:
\[
5\lvert 2x - 1 \rvert = 15 \implies \lvert 2x - 1 \rvert = 3
\]
Split into two linear equations:
\[
2x - 1 = 3 \implies 2x = 4 \implies x = 2
\]
\[
2x - 1 = -3 \implies 2x = -2 \implies x = -1
\]
Thus, the solutions are $\boxed{x = 2, x = -1}$.

(b) Sketching the graph of $y = 5\sin x - 2$ for $0^\circ \le x \le 360^\circ$\\[4pt]
* \textbf{Shape:} Sine wave over one full cycle from $0^\circ$ to $360^\circ$.
* \textbf{Amplitude:} $5$, meaning it oscillates $5$ units above and below the midline.
* \textbf{Midline:} $y = -2$.
* \textbf{Maximum value:} $-2 + 5 = 3$ at $x = 90^\circ$.
* \textbf{Minimum value:} $-2 - 5 = -7$ at $x = 270^\circ$.
* \textbf{Intercepts:} Passes through $(0^\circ, -2)$, reaches max at $(90^\circ, 3)$, crosses midline at $(180^\circ, -2)$, reaches min at $(270^\circ, -7)$, and ends at $(360^\circ, -2)$.
\end{ansbox}

%%%%%%%%%%%%%%%%%%%
\examq{4037/11/M/J/25}{6}{Differentiation}{8}
%%%%%%%%%%%%%%%%%%%
\begin{mstab}{q6color}
6(a) & Correct derivative: \newline $\frac{\mathrm{d}}{\mathrm{d}r}\left(\frac{x^2-1}{x^2+1}\right) = \frac{(x^2+1)(2x) - (x^2-1)(2x)}{(x^2+1)^2}$ \newline or $2x(x^2+1)^{-1} - 2x(x^2-1)(x^2+1)^{-2}$ & \textbf{M1} \newline \textbf{A1} & \textbf{M1} for an attempt at the quotient or product rule. \textbf{A1} fully correct. \\ \cline{2-4}
 & $\dydx = 4\left(\frac{x^2-1}{x^2+1}\right)^3 \left(\frac{(x^2+1)(2x) - (x^2-1)(2x)}{(x^2+1)^2}\right)$ \newline or $\dydx = k\left(\frac{x^2-1}{x^2+1}\right)^3 \times \mathrm{f}(x)$ & \textbf{M1} \newline \textbf{A1} & Dep \textbf{M1} for where $\mathrm{f}(x)$ is their attempt to differentiate. \textbf{A1} fully correct. \\ \cline{2-4}
 & Correct completion to $\dfrac{16x(x^2-1)^3}{(x^2+1)^5}$ & \textbf{A1} & \\ \hline
\altrow{q6color}{Alternative}
 & $u = (x^2-1)^4$, $\mathrm{d}u = 4(x^2-1)^3 \times (2x)$ oe and $v = (x^2+1)^4$, $\mathrm{d}v = 4(x^2+1)^3 \times (2x)$ oe & \textbf{(M1)} & \textbf{M1} for either correct OR for $\mathrm{d}u = 4(x^2-1)^3 \times g(x)$ where $g(x)$ is an attempt to differentiate $x^2-1$ and $\mathrm{d}v = 4(x^2+1)^3 \times h(x)$ where $h(x)$ is an attempt to differentiate $x^2+1$ \\ \hline
 & $\dydx = \dfrac{(x^2+1)^4 \times 4(x^2-1)^3(2x) - (x^2-1)^4 \times 4(x^2+1)^3(2x)}{((x^2+1)^4)^2}$ oe & \textbf{(2)} & \textbf{M1} for an attempt at the product or quotient rule \\ \hline
 & Correct completion to $\dfrac{16x(x^2-1)^3}{(x^2+1)^5}$ & \textbf{(A1)} & \\ \hline
6(b)(i) & $\text{their } 16x(x^2-1)^3 = 0$ \newline therefore $x = 0$, $x^2 - 1 = 0 \implies x = \pm 1$ oe & \textbf{B1} & \textbf{FT} $\text{their } 16$. Allow for $\text{their } 16x(x^2-1) = 0$ \\ \hline
\multirow{4}{*}{6(b)(ii)} & Any two of: \newline $-1 < x < 0$ ($\frac{\mathrm{d}y}{\mathrm{d}x}$ negative), $0$ ($0$), $0 < x < 1$ ($\frac{\mathrm{d}y}{\mathrm{d}x}$ positive) \newline $0 < x < 1$ ($\frac{\mathrm{d}y}{\mathrm{d}x}$ negative), $x > 1$ ($\frac{\mathrm{d}y}{\mathrm{d}x}$ positive), $1$ ($0$) \newline $-1 < x < 0$ ($\frac{\mathrm{d}y}{\mathrm{d}x}$ positive), $x < -1$ ($\frac{\mathrm{d}y}{\mathrm{d}x}$ negative), $-1$ ($0$) & \textbf{M1} & Must show change in the sign of first derivative around the stationary points. \textbf{M0} only if second derivative is used. \\ \cline{2-4}
 & $x = 1$ and $x = -1$ are minimum points & \textbf{A1} & \\ \hline
\end{mstab}

\begin{ansbox}{q6color}
(a) Showing the derivative $\dfrac{\mathrm{d}y}{\mathrm{d}x} = \dfrac{16x(x^2-1)^3}{(x^2+1)^5}$\\[4pt]
Using the chain rule and quotient rule on $y = \left(\frac{x^2-1}{x^2+1}\right)^4$:
\[
\dydx = 4\left(\frac{x^2-1}{x^2+1}\right)^3 \times \frac{\mathrm{d}}{\mathrm{d}x}\left(\frac{x^2-1}{x^2+1}\right)
\]
Differentiating the inner quotient:
\[
\frac{\mathrm{d}}{\mathrm{d}x}\left(\frac{x^2-1}{x^2+1}\right) = \frac{(x^2+1)(2x) - (x^2-1)(2x)}{(x^2+1)^2} = \frac{2x^3 + 2x - 2x^3 + 2x}{(x^2+1)^2} = \frac{4x}{(x^2+1)^2}
\]
Combining terms:
\[
\dydx = 4\left(\frac{(x^2-1)^3}{(x^2+1)^3}\right) \times \frac{4x}{(x^2+1)^2} = \frac{16x(x^2-1)^3}{(x^2+1)^5} \quad \text{as required.}
\]

(b)(i) Stationary points\\[4pt]
Set $\dydx = 0$:
\[
\frac{16x(x^2-1)^3}{(x^2+1)^5} = 0 \implies 16x(x^2-1)^3 = 0
\]
This yields $x = 0$ or $(x^2-1)^3 = 0 \implies x = \pm 1$. Thus, stationary points occur at $x = -1, 0, 1$.

(b)(ii) First derivative test\\[4pt]
\begin{center}
\renewcommand{\arraystretch}{1.3}
\begin{tabular}{|c|c|c|c|}
\hline
$x$ & $-1 < x < 0$ & $0$ & $0 < x < 1$ \\ \hline
$\dfrac{\mathrm{d}y}{\mathrm{d}x}$ & negative & $0$ & positive \\ \hline
\end{tabular}
\end{center}
\begin{center}
\renewcommand{\arraystretch}{1.3}
\begin{tabular}{|c|c|c|c|}
\hline
$x$ & $0 < x < 1$ & $1$ & $x > 1$ \\ \hline
$\dfrac{\mathrm{d}y}{\mathrm{d}x}$ & negative & $0$ & positive \\ \hline
\end{tabular}
\end{center}
Thus, $\boxed{x = 1 \text{ and } x = -1 \text{ are minimum points}}$.
\end{ansbox}

%%%%%%%%%%%%%%%%%%%
\examq{4037/11/M/J/25}{7}{Factors of Polynomials}{6}
%%%%%%%%%%%%%%%%%%%
\begin{mstab}{q7color}
7 & $2x^3 - 9x^2 + 10x - 3 [= 0]$ & \textbf{B1} & \\ \cline{2-4}
 & Finds a correct linear factor $x-1$ or $x-3$ or $2x-1$ & \textbf{M1} & \\ \cline{2-4}
 & Finds the correct corresponding quadratic factor: \newline For $(x-1)$ gives $2x^2 - 7x + 3$ \newline For $(x-3)$ gives $2x^2 - 3x + 1$ \newline For $(2x-1)$ gives $x^2 - 4x + 3$ & \textbf{2} & \textbf{M1} for a corresponding quadratic factor with 2 terms correct \\ \cline{2-4}
 & Factorises their 3-term quadratic factor or solves their 3-term quadratic equation & \textbf{M1} & \\ \cline{2-4}
 & $x = \frac{1}{2}, 1, 3$ & \textbf{A1} & \\ \hline
\altrow{q7color}{Alternative}
 & $2x^3 - 9x^2 + 10x - 3 = 0$ & \textbf{(B1)} & \\ \hline
 & Finds a correct linear factor $x-1$ or $x-3$ or $2x-1$ & \textbf{(M1)} & \\ \hline
 & Finds a second correct linear factor $x-1$ or $x-3$ or $2x-1$ & \textbf{(M1)} & \\ \hline
 & Finds a third correct linear factor $x-1$ or $x-3$ or $2x-1$ & \textbf{(M1)} & \\ \hline
 & $x-1$, $x-3$ and $2x-1$ & \textbf{(A1)} & \\ \hline
 & $x = \frac{1}{2}, 1, 3$ & \textbf{(A1)} & \\ \hline
\end{mstab}

\begin{ansbox}{q7color}
Finding the $x$-coordinates where the curve cuts the $x$-axis\\[4pt]
Set $y = 0$ in the curve equation $y = (2x-9)(x^2+5) + 42$:
\[
(2x-9)(x^2+5) + 42 = 0
\]
Expand the brackets:
\[
2x^3 + 10x - 9x^2 - 45 + 42 = 0 \implies 2x^3 - 9x^2 + 10x - 3 = 0
\]
Let $\mathrm{p}(x) = 2x^3 - 9x^2 + 10x - 3$. Test $x = 1$:
\[
\mathrm{p}(1) = 2(1)^3 - 9(1)^2 + 10(1) - 3 = 2 - 9 + 10 - 3 = 0
\]
Thus, $(x-1)$ is a factor. Perform polynomial long division or equating coefficients to find the quadratic factor:
\[
2x^3 - 9x^2 + 10x - 3 = (x-1)(2x^2 - 7x + 3) = 0
\]
Factorise the quadratic expression $2x^2 - 7x + 3$:
\[
2x^2 - 6x - x + 3 = 2x(x-3) - 1(x-3) = (2x-1)(x-3)
\]
So the complete factorization is:
\[
(x-1)(2x-1)(x-3) = 0
\]
Solving for $x$, we get $\boxed{x = 1, x = \frac{1}{2}, x = 3}$.
\end{ansbox}

%%%%%%%%%%%%%%%%%%%
\examq{4037/11/M/J/25}{8}{Logarithmic \& Exponential Functions}{7}
%%%%%%%%%%%%%%%%%%%
\begin{mstab}{q8color}
8(a) & $x > \frac{4}{12}$ oe & \textbf{B1} & allow $0.3$ for $\frac{4}{12}$ \\ \hline
8(b) & $\log_x 125 = \frac{\log_5 125}{\log_5 x} = \frac{3}{\log_5 x}$ \newline or $\frac{1}{\log_x 125} = \log_{125} x$ \newline or $\log_5(12x-4) = \frac{\log_x(12x-4)}{\log_x 5}$ & \textbf{B1} & for a correct and relevant change of base seen \\ \cline{2-4}
 & $\log_5(12x-4) = 2\log_5 x + 1$ \newline or $\log_x(12x-4) = 2 + \log_x 5$ & \textbf{B1} & Dep on correct change of base \\ \cline{2-4}
 & $\log_5\left(\frac{12x-4}{x^2}\right) = 1$ oe \newline or $\log_x\left(\frac{12x-4}{5}\right) = 2$ oe \newline or $\log_5(12x-4) = \log_5(5x^2)$ & \textbf{B1} & Dep on correct change of base \\ \cline{2-4}
 & $5x^2 - 12x + 4 = 0$ & \textbf{B1} & Dep on correct change of base \\ \cline{2-4}
 & $(5x-2)(x-2) = 0$ oe & \textbf{M1} & dep on at least one previous \textbf{B1} awarded \\ \cline{2-4}
 & $x = 0.4, 2$ & \textbf{A1} & \\ \hline
\end{mstab}

\begin{ansbox}{q8color}
(a) Range of values for which $\log_5(12x-4)$ exists\\[4pt]
The argument of a logarithm must be strictly positive:
\[
12x - 4 > 0 \implies 12x > 4 \implies x > \frac{4}{12} \implies \boxed{x > \frac{1}{3} \text{ (or } x > 0.333\text{)}}
\]

(b) Solving the logarithmic equation\\[4pt]
Given equation:
\[
\log_5(12x-4) = \frac{6}{\log_x 125} + 1
\]
Using the change of base rule, $\log_x 125 = \frac{\log_5 125}{\log_5 x} = \frac{3}{\log_5 x}$. Thus:
\[
\frac{6}{\log_x 125} = \frac{6}{\frac{3}{\log_5 x}} = 2\log_5 x
\]
Substitute back into the equation:
\[
\log_5(12x-4) = 2\log_5 x + 1
\]
Rewrite $1$ as $\log_5 5$ and $2\log_5 x$ as $\log_5(x^2)$:
\[
\log_5(12x-4) = \log_5(x^2) + \log_5 5 = \log_5(5x^2)
\]
Equating arguments:
\[
12x - 4 = 5x^2 \implies 5x^2 - 12x + 4 = 0
\]
Factorising the quadratic:
\[
(5x - 2)(x - 2) = 0 \implies x = \frac{2}{5} = 0.4 \text{ or } x = 2
\]
Both values satisfy $x > \frac{1}{3}$. Thus, $\boxed{x = 0.4, 2}$.
\end{ansbox}

%%%%%%%%%%%%%%%%%%%
\examq{4037/11/M/J/25}{9}{Differentiation \textperiodcentered\ Integration}{10}
%%%%%%%%%%%%%%%%%%%
\begin{mstab}{q9color}
9 & $\left[\text{When } x = 2\right] y = 3$ & \textbf{B1} & \\ \cline{2-4}
 & $\dydx = \frac{1}{2}(4x+1)^{-\frac{1}{2}} \times 4$ & \textbf{M1} & for $\dydx = \frac{1}{2}(4x+1)^{-\frac{1}{2}} \times k$, $k \neq 0$ \\ \cline{2-4}
 & $\left[\text{gradient of tangent} = \right] \frac{2}{3}$ & \textbf{M1} & Dep on first \textbf{M1} \\ \cline{2-4}
 & $y - \text{their } 3 = \frac{2}{3}(x-2)$ oe & \textbf{A1} & Must be from correct straight line equation \\ \cline{2-4}
 & $\left[x\text{-intercept for tangent} = \right] 1 - \frac{5}{2}$ soi $\left(= -\frac{1}{2}\right)$ or similar & \textbf{B1} & \\ \cline{2-4}
 & $\frac{1}{2} \times \left(\frac{5}{2} + 2\right) \times 3$ soi & \textbf{B1} & For area of triangle allow unsimplified \\ \cline{2-4}
 & $\left[x\text{-intercept for curve} = \right] -\frac{1}{4}$ soi & \textbf{B1} & \\ \cline{2-4}
 & Area below curve: $\displaystyle\int_{-\frac{1}{4}}^{2} (4x+1)^{\frac{1}{2}} \,\dd x$ & \textbf{M1} & \textbf{M1} for integration in the form of $k(4x+1)^{\frac{3}{2}}$, $k \neq 0$ \\ \cline{2-4}
 & $= \left[ \dfrac{(4x+1)^{\frac{3}{2}}}{\frac{3}{2} \times 4} \right]_{-\frac{1}{4}}^{2} = \left[ \frac{(4x+1)^{\frac{3}{2}}}{6} \right]_{-\frac{1}{4}}^{2}$ oe & \textbf{M1} & Dep on previous \textbf{M1} for area under the curve \\ \cline{2-4}
 & Correct use of upper and lower limits: $\left(\frac{(4(2)+1)^{\frac{3}{2}}}{6} - \frac{(4(-\frac{1}{4})+1)^{\frac{3}{2}}}{6}\right)$ oe & \textbf{A1} & Allow if using wrong limits but must see correct substitution \\ \cline{2-4}
 & $= \frac{27}{4} - \frac{9}{2} = \frac{9}{4}$ oe & \textbf{M1} & \\ \cline{2-4}
 & Shaded area $= \text{Area of triangle} - \text{Area under curve} = \frac{15}{4} - \frac{9}{4} = \frac{6}{4} = \frac{3}{2}$ & \textbf{A1} & \\ \hline
\end{mstab}

\begin{ansbox}{q9color}
Finding the area of the shaded region\\[4pt]
The curve is given by $y = \sqrt{4x+1}$. Point $A$ has $x$-coordinate $2$, so its $y$-coordinate is $y = \sqrt{4(2)+1} = \sqrt{9} = 3$. 

\textbf{1. Equation of the tangent at $A$:}
Find the derivative $\dydx$:
\[
\dydx = \frac{1}{2}(4x+1)^{-\frac{1}{2}} \times 4 = \frac{2}{\sqrt{4x+1}}
\]
At $x = 2$, the gradient of the tangent is:
\[
m = \frac{2}{\sqrt{4(2)+1}} = \frac{2}{3}
\]
The tangent equation passing through $A(2,3)$ is:
\[
y - 3 = \frac{2}{3}(x - 2) \implies y = \frac{2}{3}x + \frac{5}{3}
\]
The $x$-intercept of this tangent occurs when $y = 0$:
\[
0 = \frac{2}{3}x + \frac{5}{3} \implies 2x + 5 = 0 \implies x = -\frac{5}{2}
\]

\textbf{2. Area of the triangle above the $x$-axis:}
The tangent forms a triangle with the $x$-axis from $x = -\frac{5}{2}$ to $x = 2$, with a height of $3$ at $x = 2$:
\[
\text{Area}_{\text{triangle}} = \frac{1}{2} \times \text{base} \times \text{height} = \frac{1}{2} \times \left(2 - \left(-\frac{5}{2}\right)\right) \times 3 = \frac{1}{2} \times \frac{9}{2} \times 3 = \frac{27}{4}
\]

\textbf{3. Area under the curve:}
The curve cuts the $x$-axis when $y = 0 \implies 4x + 1 = 0 \implies x = -\frac{1}{4}$. Integrate the curve from $x = -\frac{1}{4}$ to $x = 2$:
\[
\text{Area}_{\text{curve}} = \int_{-\frac{1}{4}}^{2} (4x+1)^{\frac{1}{2}} \,\dd x = \left[ \frac{(4x+1)^{\frac{3}{2}}}{\frac{3}{2} \times 4} \right]_{-\frac{1}{4}}^{2} = \left[ \frac{(4x+1)^{\frac{3}{2}}}{6} \right]_{-\frac{1}{4}}^{2}
\]
Substitute the limits:
\[
\text{Area}_{\text{curve}} = \frac{(9)^{\frac{3}{2}}}{6} - \frac{(0)^{\frac{3}{2}}}{6} = \frac{27}{6} - 0 = \frac{9}{2}
\]

\textbf{4. Shaded area:}
\[
\text{Shaded Area} = \text{Area}_{\text{triangle}} - \text{Area}_{\text{curve}} = \frac{27}{4} - \frac{9}{2} = \frac{27}{4} - \frac{18}{4} = \boxed{\frac{9}{4} \text{ (or } 2.25\text{)}}
\]
\end{ansbox}

%%%%%%%%%%%%%%%%%%%
\examq{4037/11/M/J/25}{10}{Trigonometry}{7}
%%%%%%%%%%%%%%%%%%%
\begin{mstab}{q1color}
10(a) & $\dfrac{\sin\theta}{\sqrt{\cot^2\theta}} + \dfrac{1}{\sqrt{\sec^2\theta}}$ & \textbf{B1} & for correct use of the trigonometry identity $\operatorname{cosec}^2\theta - 1 = \cot^2\theta$ \\ \cline{2-4}
 & $\sin\theta\tan\theta + \cos\theta$ oe & \textbf{B1} & \\ \cline{2-4}
 & $\sin\theta \times \frac{\sin\theta}{\cos\theta} + \cos\theta$ & \textbf{B1} & \\ \cline{2-4}
 & $\frac{\sin^2\theta + \cos^2\theta}{\cos\theta} = \frac{1}{\cos\theta} = \sec\theta$ & \textbf{B1} & Dep on all 3 previous marks. Last mark is not available if persistent missing of $\theta$. \\ \hline
\altrow{q1color}{Alternative}
 & $\dfrac{\sin\theta}{\sqrt{\cot^2\theta}} + \dfrac{1}{\sqrt{\sec^2\theta}}$ & \textbf{(B1)} & for correct use of the trigonometry identity $\operatorname{cosec}^2\theta - 1 = \cot^2\theta$ \\ \hline
 & $\sec\theta\sin\theta + \cot\theta$ & \textbf{(B1)} & \\ \hline
 & $\frac{\sin\theta}{\cos\theta} + \frac{\cos\theta}{\sin\theta}\sec\theta \dots$ & \textbf{(B1)} & \\ \hline
 & $\frac{\sin^2\theta + \cos^2\theta}{\sin\theta\cos\theta} \dots$ & \textbf{(B1)} & Dep on all 3 previous marks \\ \hline
10(b) & $\cos x = \frac{1}{\alpha}$ & \textbf{B1} & $\frac{1}{\alpha^2} = 1 - \sin^2 x$ could be implied by \\ \cline{2-4}
 & $\sin^2 x + \left(\frac{1}{\alpha}\right)^2 = 1$ oe OR $[\text{opposite}^2 = ] \alpha^2 - 1$ soi and use of $\sin x = \frac{\text{opposite}}{\text{hypotenuse}}$ at some stage leading to $\sin x = [\pm]\sqrt{1 - \frac{1}{\alpha^2}}$ oe & \textbf{M1} & \\ \cline{2-4}
 & $\sin x = -\sqrt{1 - \frac{1}{\alpha^2}}$ oe & \textbf{A1} & \\ \hline
\end{mstab}

\begin{ansbox}{q1color}
(a) Proving the trigonometric identity\\[4pt]
Start with the left-hand side and use the identity $\operatorname{cosec}^2\theta - 1 = \cot^2\theta$ and $\sec^2\theta = 1 + \tan^2\theta$:
\[
\text{LHS} = \frac{\sin\theta}{\sqrt{\operatorname{cosec}^2\theta - 1}} + \frac{1}{\sqrt{\sec^2\theta}} = \frac{\sin\theta}{\sqrt{\cot^2\theta}} + \frac{1}{\sec\theta}
\]
Since $\theta$ is acute ($0 \le \theta < \frac{\pi}{2}$), $\sqrt{\cot^2\theta} = \cot\theta$ and $\frac{1}{\sec\theta} = \cos\theta$:
\[
= \frac{\sin\theta}{\frac{\cos\theta}{\sin\theta}} + \cos\theta = \frac{\sin^2\theta}{\cos\theta} + \cos\theta
\]
Put over a common denominator:
\[
= \frac{\sin^2\theta + \cos^2\theta}{\cos\theta}
\]
Since $\sin^2\theta + \cos^2\theta = 1$:
\[
= \frac{1}{\cos\theta} = \sec\theta = \text{RHS} \quad \text{as required.}
\]

(b) Finding $\sin x$ in terms of $\alpha$\\[4pt]
Given $\sec x = \alpha$, which means $\cos x = \frac{1}{\alpha}$. We are given the domain $\frac{3\pi}{2} < x \le 2\pi$, which is the fourth quadrant where sine is negative and cosine is positive.
Using the Pythagorean identity $\sin^2 x + \cos^2 x = 1$:
\[
\sin^2 x + \left(\frac{1}{\alpha}\right)^2 = 1 \implies \sin^2 x = 1 - \frac{1}{\alpha^2} = \frac{\alpha^2 - 1}{\alpha^2}
\]
Taking square roots and choosing the negative sign for the fourth quadrant:
\[
\sin x = -\sqrt{1 - \frac{1}{\alpha^2}} \text{ or } \boxed{-\frac{\sqrt{\alpha^2-1}}{\alpha}}
\]
\end{ansbox}

%%%%%%%%%%%%%%%%%%%
\examq{4037/11/M/J/25}{11}{Arithmetic \& Geometric Progressions}{9}
%%%%%%%%%%%%%%%%%%%
\begin{mstab}{q2color}
11(a) & $10 - 2d$ and $10 - d$ & \textbf{B2} & \textbf{B1} for either, ignore labels \\ \hline
11(b) & Squares each term and forms a sum e.g. $(10-2d)^2 + (10-d)^2 + 10^2 = 140$ or $a^2 + (a+d)^2 + (a+2d)^2 = 140$ & \textbf{M1} & \textbf{FT} from their two terms in (a), must be both 2 terms expression in terms of $d$ only or correct expressions in terms of $a$ only, must see substitution and expansion \\ \cline{2-4}
 & Correct equation in solvable form: $5d^2 - 60d + 160 = 0$ oe or $5a^2 + 20a - 60 = 0$ oe & \textbf{A1} & \\ \cline{2-4}
 & Factorises their 3-term quadratic expression or solves their 3-term quadratic equation & \textbf{M1} & \\ \cline{2-4}
 & $d = 8, d = 4$ & \textbf{A1} & \\ \cline{2-4}
 & Finds the first term when $d = 4$: $a = 2$ & \textbf{M1} & \textbf{FT} $10 - 2(\text{their } d)$ \\ \cline{2-4}
 & $\left[S_{200} = \right] \frac{200}{2} \{4 + 199 \times 4\}$ oe & \textbf{M1} & \textbf{FT} $\text{their } a$ and $\text{their } d$ \\ \cline{2-4}
 & $80\,000$ & \textbf{A1} & \\ \hline
\end{mstab}

\begin{ansbox}{q2color}
(a) Expressions for the 1st and 2nd terms in terms of $d$\\[4pt]
Let the arithmetic progression have 3rd term $u_3 = 10$ and common difference $d$. Using $u_n = a + (n-1)d$:
* 3rd term: $u_3 = a + 2d = 10 \implies a = 10 - 2d$
* 1st term ($a$): $\boxed{10 - 2d}$
* 2nd term ($u_2 = a + d$): $(10 - 2d) + d = \boxed{10 - d}$

(b) Sum of the first 200 terms for the smaller value of $d$\\[4pt]
When each of the first 3 terms is squared, their sum is $140$:
\[
(10-2d)^2 + (10-d)^2 + 10^2 = 140
\]
Expand each term:
\[
(100 - 40d + 4d^2) + (100 - 20d + d^2) + 100 = 140
\]
Combine like terms:
\[
5d^2 - 60d + 300 = 140 \implies 5d^2 - 60d + 160 = 0
\]
Divide by $5$:
\[
d^2 - 12d + 32 = 0 \implies (d - 4)(d - 8) = 0 \implies d = 4 \text{ or } d = 8
\]
The smaller value is $d = 4$. Find the first term $a$ using $d = 4$:
\[
a = 10 - 2(4) = 10 - 8 = 2
\]
Now find the sum of the first 200 terms ($S_{200}$) using $S_n = \frac{n}{2}[2a + (n-1)d]$:
\[
S_{200} = \frac{200}{2} [2(2) + (200 - 1)(4)] = 100 [4 + 199 \times 4] = 100 [4 + 796] = 100 \times 800 = \boxed{80\,000}
\]
\end{ansbox}

%%%%%%%%%%%%%%%%%%%
\examq{4037/11/M/J/25}{12}{Permutations \& Combinations}{4}
%%%%%%%%%%%%%%%%%%%
\begin{mstab}{q3color}
12 & For substituting $\ncr{n}{5}$ and $\ncr{n-1}{5}$ e.g. $\frac{n!}{(n-5)!5!} - \frac{(n-1)!}{(n-6)!5!}$ & \textbf{M1} & \\ \cline{2-4}
 & For factorising $\frac{(n-1)!}{(n-6)!5!}$ or $\frac{(n-1)!}{(n-5)!5!}$ e.g. $\frac{(n-1)!}{(n-6)!5!} \left(\frac{n}{n-5} - 1\right)$ oe & \textbf{M1} & Dep on previous \textbf{M1} awarded. Do not allow if using $^{n-1}\mathrm{C}_4$ to simplify the left-hand side. \\ \cline{2-4}
 & Simplifying fraction: $\frac{(n-1)!}{(n-6)!5!} \left(\frac{5}{n-5}\right)$ oe & \textbf{M1} & Dep on at least one previous \textbf{M1} awarded \\ \cline{2-4}
 & Completes argument: $\frac{(n-1)!}{(n-6)!5!} \left(\frac{5}{n-5}\right) = \frac{(n-1)!}{(n-1-4)!4!} = {}^{n-1}\mathrm{C}_4$ & \textbf{A1} & \\ \hline
\altrow{q3color}{Alternative 1}
 & For substituting $\ncr{n}{5}$ and $\ncr{n-1}{5}$ & \textbf{(M1)} & \\ \hline
 & $\frac{1}{(n-6)!} = \frac{n-5}{(n-5)!}$ or $n! = n \times (n-1)!$ & \textbf{(B1)} & Dep on previous \textbf{M1} awarded \\ \hline
 & Factorising $(n-1)!$ in the fraction: $\frac{(n-1)! \times [n - (n-5)]}{(n-5)!5!}$ & \textbf{(M1)} & Dep on previous \textbf{M1} awarded \\ \hline
 & $\frac{(n-1)! \times 5}{(n-5)!5!} = \frac{(n-1)!}{(n-5)!4!} = {}^{n-1}\mathrm{C}_4$ & \textbf{(A1)} & \\ \hline
\altrow{q3color}{Alternative 2}
 & For substituting $\ncr{n}{5}$ and $\ncr{n-1}{5}$ could be implied by expansions of the factorial & \textbf{(M1)} & \\ \hline
 & $\frac{n(n-1)(n-2)(n-3)(n-4)}{5!} - \frac{(n-1)(n-2)(n-3)(n-4)(n-5)}{5!}$ & \textbf{(B1)} & Dep on previous \textbf{M1} awarded \\ \hline
 & Factorising $(n-1)(n-2)(n-3)(n-4)$ in the fraction: $\frac{(n-1)(n-2)(n-3)(n-4)[n - (n-5)]}{5!}$ & \textbf{(M1)} & Dep on previous \textbf{M1} awarded \\ \hline
 & Completes argument: $\frac{(n-1)(n-2)(n-3)(n-4)}{4!} = \frac{(n-1)!}{4!(n-5)!} = {}^{n-1}\mathrm{C}_4$ & \textbf{(A1)} & \\ \hline
\end{mstab}

\begin{ansbox}{q3color}
Proving ${}^n\mathrm{C}_5 - {}^{n-1}\mathrm{C}_5 = {}^{n-1}\mathrm{C}_4$\\[4pt]
Substitute the definition of combinations $\ncr{n}{r} = \frac{n!}{(n-r)!r!}$ into the left-hand side:
\[
\text{LHS} = {}^n\mathrm{C}_5 - {}^{n-1}\mathrm{C}_5 = \frac{n!}{(n-5)!5!} - \frac{(n-1)!}{((n-1)-5)!5!} = \frac{n!}{(n-5)!5!} - \frac{(n-1)!}{(n-6)!5!}
\]
Express $n!$ as $n \times (n-1)!$ and $(n-5)!$ as $(n-5) \times (n-6)!$:
\[
= \frac{n(n-1)!}{(n-5)(n-6)!5!} - \frac{(n-1)!}{(n-6)!5!}
\]
Factor out the common term $\frac{(n-1)!}{(n-6)!5!}$:
\[
= \frac{(n-1)!}{(n-6)!5!} \left[ \frac{n}{n-5} - 1 \right]
\]
Simplify the expression inside the brackets:
\[
\frac{n}{n-5} - 1 = \frac{n - (n-5)}{n-5} = \frac{5}{n-5}
\]
Combine terms:
\[
= \frac{(n-1)!}{(n-6)!5!} \times \frac{5}{n-5} = \frac{(n-1)!}{(n-6)!(n-5) \times \frac{5!}{5}} = \frac{(n-1)!}{(n-5)! \times 4!}
\]
Recognise this as the standard formula for combinations:
\[
= \frac{(n-1)!}{((n-1)-4)!4!} = {}^{n-1}\mathrm{C}_4 = \text{RHS} \quad \text{as required.}
\]
\end{ansbox}

\end{document}